\documentclass[11pt]{article}
\usepackage[margin=1.1in]{geometry}
\usepackage{amsmath,amssymb,amsthm}
\usepackage{booktabs}
\usepackage[hidelinks]{hyperref}
\usepackage{microtype}

\newtheorem{deriv}{Derivation}
\theoremstyle{remark}
\newtheorem*{scope}{Scope}
\newcommand{\Nm}{\mathrm{N}}
\newcommand{\bucket}[1]{\par\smallskip\noindent\textit{Status: #1.}}
\newcommand{\chk}[1]{\par\noindent\textit{Verified numerically:
\texttt{#1} in \texttt{out/derivations\_check.json}.}}

\title{Derivations for the Prime Pages:\\
\large every formula used by the explainer and the explorables ---
derived in full, defined with source, or named as a gap}
\author{%
  Lee Smart\thanks{Companion to the explorables
(\texttt{explorables/}), the five-pictures explainer
(\texttt{papers/how-the-primes-work.pdf}) and the ledger
(\texttt{papers/prime-phenomena-ledger.pdf}) in
\texttt{github.com/vfd-org/closure-positivity-lab}. Selected numeric
consequences of the derivations (18 checks, named per derivation in
the text) are verified by
\texttt{python3 -m lab.derivations\_check}; the machine-readable
results ship at \texttt{out/derivations\_check.json}.}\\[4pt]
  \textit{Institute of Vibrational Field Dynamics}\\[2pt]
  \texttt{contact@vibrationalfielddynamics.org}\\[2pt]
  \texttt{@vfd\_org}%
}
\date{2026-06-12}

\begin{document}
\maketitle

\begin{abstract}\noindent
The prime pages of this repository use a small set of formulas: wheel
survival counts, the singular-series factors, the Hardy--Littlewood
count, Sato--Tate moments and their gate variances, the Poisson and GUE
baselines, the Brandt trace-extraction identities, the Eichler mass, the
tapered prime spectrum, the Fibonacci-chain peak module, and the
Chebyshev bias mechanism. This note derives each one at the level it
deserves, under a strict three-bucket rule: \emph{derived in full} (the
derivation is complete on the page), \emph{defined with source} (a deep
theorem used as input, never re-proven and never assumed beyond its
statement), or \emph{named as a gap} (heuristic or open --- said so).
Nothing receives a fourth category. Eighteen numeric checks verify
selected consequences, named per derivation in the text.
\end{abstract}

\begin{scope}
All mathematics here is classical; the derivations are expositions of
known material, organised for the explorables, and we claim no priority
for any of them. Heuristic steps are labelled at the point of use.
Nothing in this note proves or claims progress on any open conjecture.
\end{scope}

\section{Part I: the wheels and the singular series}

\begin{deriv}[D1: wheel survival count]\label{d1}
Fix a pattern of $k$ distinct integer offsets
$H = \{h_1, \dots, h_k\}$ and a prime $p$. A residue $r \bmod p$
\emph{survives} the wheel of $p$ if $r + h_i \not\equiv 0 \pmod p$ for
every $i$. The forbidden residues are exactly
$\{-h_i \bmod p\}$, of which there are $w(p) := \#\{h_i \bmod p\}$
distinct values (offsets congruent mod $p$ forbid the same cell). Hence
exactly $p - w(p)$ residues survive, and the survival fraction is
$1 - w(p)/p$. If $w(p) = p$ the pattern is \emph{inadmissible}: every
cell dies, and only finitely many occurrences can exist (each must
include the prime $p$ itself, as in $(n, n{+}2, n{+}4)$ at $p = 3$,
realised once by $(3,5,7)$).
\bucket{derived in full}
\end{deriv}

\begin{deriv}[D2: the independence baseline]\label{d2}
The comparison model treats the $k$ shifted values as unrelated
integers. An integer avoids $0 \bmod p$ in $p-1$ of $p$ residues, so
$k$ independent integers all avoid it with probability
$\left((p-1)/p\right)^k$. This is a \emph{model}, not a theorem about
primes; it is the denominator against which the wheels' joint behaviour
is measured.
\bucket{derived in full (as a definition of the baseline)}
\end{deriv}

\begin{deriv}[D3: the per-wheel factor]\label{d3}
The wheel of $p$ corrects the independence baseline by
\[
\delta_p(H) \;=\; \frac{1 - w(p)/p}{(1 - 1/p)^k},
\]
the ratio of the actual joint survival fraction (D1) to the independent
one (D2). For all $p$ exceeding every pairwise difference $|h_i - h_j|$
the offsets are distinct mod $p$, so $w(p) = k$.
\bucket{derived in full}
\end{deriv}

\begin{deriv}[D4: the gap-ratio law]\label{d4}
Specialise to pairs, $H = \{0, g\}$ with $g$ even, and an odd prime
$p$. If $p \nmid g$ then $w(p) = 2$ and
$\delta_p = (1 - 2/p)(1 - 1/p)^{-2} = p(p-2)/(p-1)^2$ --- the twin
factor. If $p \mid g$ then $0 \equiv g$, so $w(p) = 1$ and
$\delta_p = (1 - 1/p)(1-1/p)^{-2} = p/(p-1)$. The enhancement of a gap
$g$ over the twin gap, per odd prime $p \mid g$, is therefore
\[
\frac{p/(p-1)}{\,p(p-2)/(p-1)^2\,} \;=\; \frac{p-1}{p-2},
\qquad\text{so}\qquad
w(g) \;=\; \prod_{\substack{p \mid g \\ p \text{ odd}}} \frac{p-1}{p-2}.
\]
Hence $w(6) = 2$ (from $p{=}3$), $w(10) = 4/3$ (from $p{=}5$), and
$w(210) = 2 \cdot \tfrac43 \cdot \tfrac65 = \tfrac{16}{5}$ (from
$3, 5, 7$). These are the gap-slider's predictions, and the sieve
measures $1.9950$, $1.3318$, $3.1994$ at $5\times10^7$ (ledger row 2).
\bucket{derived in full}\chk{D4}
\end{deriv}

\begin{deriv}[D5: the two forms of the twin constant]\label{d5}
$\dfrac{p(p-2)}{(p-1)^2} = \dfrac{(p-1)^2 - 1}{(p-1)^2}
= 1 - \dfrac{1}{(p-1)^2}$, so
\[
C_2 \;=\; \prod_{p > 2} \frac{p(p-2)}{(p-1)^2}
      \;=\; \prod_{p > 2} \Bigl(1 - \frac{1}{(p-1)^2}\Bigr)
      \;=\; 0.6601618158\ldots
\]
The first form is the wheel product (D4); the second is the standard
Hardy--Littlewood form. They are the same number by one line of algebra.
\bucket{derived in full}\chk{D5}
\end{deriv}

\begin{deriv}[D6: convergence of the singular series]\label{d6}
For $p$ beyond the offsets' span, $w(p) = k$ and
\[
\log \delta_p = \log\Bigl(1 - \frac{k}{p}\Bigr)
              - k \log\Bigl(1 - \frac{1}{p}\Bigr)
= -\frac{k(k-1)}{2p^2} + \frac{k - k^3}{3p^3} + O(p^{-4}),
\]
by expanding both logarithms (the $1/p$ terms cancel exactly --- this
cancellation is the reason the product converges). Since
$\sum_p p^{-2} < \infty$, the product
$\mathfrak{S}(H) = \prod_p \delta_p(H)$ converges absolutely for every
admissible pattern. For twins ($k = 2$) the second-order coefficient is
$-1$, matching $1 - 1/(p-1)^2 = 1 - p^{-2} + O(p^{-3})$.
\bucket{derived in full}\chk{D6, both expansion coefficients}
\end{deriv}

\begin{deriv}[D7: the Hardy--Littlewood count, and the lower
limit]\label{d7}
\emph{Heuristic step, labelled as such.} In the Cram\'er model an
integer near $t$ is prime with probability $1/\log t$; $k$ independent
such events give $1/\log^k t$; the wheels correct the independence
assumption place-by-place, multiplying by $\mathfrak{S}(H)$ (D3, D6).
Summing over starting points,
\[
\#\{n \le x : n + h_i \text{ all prime}\}
\;\approx\;
\mathfrak{S}(H) \int_2^x \frac{dt}{\log^k t}.
\]
This \emph{is} the Hardy--Littlewood $k$-tuple conjecture: the wheels
prove admissibility and the constant; the asymptotic itself is open for
every $k \ge 2$ (gap 1 of the mathematics page). Two facts about the
integral that the explorables rely on: (i) any fixed lower limit
$A \ge 2$ changes the integral by the constant
$\int_2^A dt/\log^k t$, so all choices are asymptotically equivalent;
(ii) at finite $x$ and large $k$ that constant is \emph{not} small
relative to the total (at $k = 6$, $x = 2\times10^7$ the $[2, 100]$
portion is over half), which is why the wheel-lab integrates from
$100$ and says so on the page.
\bucket{heuristic, named as such; the asymptotic is the open
conjecture}
\end{deriv}

\begin{deriv}[D8: the Goldbach factor and the leading 2]\label{d8}
Representations $N = a + (N - a)$ with both parts prime. At an odd
prime $p \nmid N$: $a$ must avoid $\{0, N \bmod p\}$, two distinct
cells, giving the twin factor $p(p-2)/(p-1)^2$. At an odd $p \mid N$:
the two forbidden cells coincide, giving $p/(p-1)$ --- the same
per-prime enhancement $(p-1)/(p-2)$ as in D4. At $p = 2$ with $N$
even: $a$ must be odd, survival $1/2$, against baseline
$(1/2)^2 = 1/4$ --- a factor of exactly $2$. Collecting:
\[
R(N) \;\approx\; 2\,C_2 \prod_{\substack{p \mid N\\ p \text{ odd}}}
\frac{p-1}{p-2} \int_2^{N-2} \frac{dt}{\log t\,\log(N-t)},
\]
with the same heuristic status as D7. The leading $2$ in $2C_2$ ---
in the twin count as well --- is the $p = 2$ wheel: even patterns
survive the parity wheel twice as often as independence predicts.
\bucket{derived in full (the factor structure; the count itself has D7's heuristic status)}
\chk{D8}
\end{deriv}

\section{Part II: the instrument's gates}

\begin{deriv}[D9: Sato--Tate moments and the gate variances]\label{d9}
For the measure $d\mu = \frac{2}{\pi}\sin^2\theta\,d\theta$ on
$[0, \pi]$ and $x = \cos\theta$: using
$\cos^{2m}\theta\,\sin^2\theta = \cos^{2m}\theta - \cos^{2m+2}\theta$
and $\int_0^\pi \cos^{2j}\theta\,d\theta = \pi\binom{2j}{j}4^{-j}$,
\[
I_{2m} := \int x^{2m}\,d\mu
= 2\left[\binom{2m}{m}4^{-m} - \binom{2m+2}{m+1}4^{-(m+1)}\right]
= \frac{\binom{2m}{m}}{(m+1)\,4^m}
= \frac{\mathrm{Cat}(m)}{4^m},
\]
since $\binom{2m+2}{m+1}/\binom{2m}{m} = 2(2m+1)/(m+1)$ makes the
bracket collapse to $\binom{2m}{m}4^{-m}/(2(m+1))$. So
$I_2 = \tfrac14$, $I_4 = \tfrac18$, $I_6 = \tfrac{5}{64}$,
$I_8 = \tfrac{14}{256}$; odd moments vanish by the symmetry
$\theta \mapsto \pi - \theta$. The ledger's $3\sigma$ gates use the
null variances of the empirical moments over $n$ angles:
$\mathrm{Var}(\overline{x^j}) = (I_{2j} - I_j^2)/n$, i.e.
\[
\sigma(m_1) = \tfrac{1}{2\sqrt n},\quad
\sigma(m_2) = \tfrac{1}{4\sqrt n},\quad
\sigma(m_3) = \sqrt{\tfrac{5}{64 n}},\quad
\sigma(m_4) = \sqrt{\tfrac{10}{256 n}},
\]
which are exactly the constants used by
\texttt{row07\_sato\_tate.py}.
\bucket{derived in full}\chk{D9 (six checks)}
\end{deriv}

\begin{deriv}[D10: the Poisson and GUE baselines]\label{d10}
Poisson (uncorrelated) unit-mean spacings are exponential:
$P(s < 0.2) = 1 - e^{-0.2} = 0.18127\ldots$, the ``$0.18$'' of ledger
row 10. The GUE Wigner surmise is
$p(s) = \frac{32}{\pi^2}s^2 e^{-4s^2/\pi}$; with $a = 4/\pi$ and the
Gaussian moments $\int_0^\infty s^3 e^{-as^2} ds = \frac{1}{2a^2}$,
$\int_0^\infty s^4 e^{-as^2} ds = \frac{3}{8}\sqrt{\pi}\,a^{-5/2}$:
\[
\mathbb{E}[s] = \frac{32}{\pi^2}\cdot\frac{\pi^2}{32} = 1,
\qquad
\mathbb{E}[s^2] = \frac{32}{\pi^2}\cdot\frac{3\sqrt{\pi}}{8}
\Bigl(\frac{\pi}{4}\Bigr)^{5/2} = \frac{3\pi}{8},
\]
so the surmise variance is $3\pi/8 - 1 = 0.17810\ldots$ --- the
``$0.178$'' gate. The surmise itself is an approximation to the exact
GUE spacing law (defined with source: Mehta); the ledger gates against
the surmise with the documented finite-height allowance.
\bucket{derived in full (from the stated surmise density)}
\chk{D10 (three checks)}
\end{deriv}

\begin{deriv}[D11: Eisenstein eigenvalue and trace extraction]\label{d11}
The Brandt matrix $B(P)$ at a good prime $P$ has constant row sums
$\Nm(P) + 1 = \#\mathbb{P}^1(\mathcal{O}_K/P)$ (each ideal class
receives all $\Nm(P)+1$ subideals of norm $\Nm(P)$, weighted), so the
all-ones vector is an eigenvector with eigenvalue $\Nm(P)+1$: the
Eisenstein class. With Brandt dimension $h = 2$ the remaining
eigenvalue is read off the trace,
$a_P = \operatorname{tr} B(P) - (\Nm(P) + 1)$. With $h = 3$ and a
genus-2 cuspidal block whose eigenvalues are a Galois pair
$\alpha, \sigma(\alpha) \in \mathbb{Z}[\varphi]$:
\[
\operatorname{Tr}(\alpha) = \operatorname{tr} B - (\Nm(P)+1),
\qquad
\mathrm{N}(\alpha) = \frac{\det B}{\Nm(P)+1},
\]
since trace and determinant collect eigenvalues additively and
multiplicatively. The test for real multiplication: $\alpha$ generates
$\mathbb{Q}(\sqrt{\Delta})$ with
$\Delta = \operatorname{Tr}^2 - 4\mathrm{N}$, so
$\alpha \in \mathbb{Z}[\varphi] \setminus \mathbb{Z}$ iff
$\Delta = 5 \times (\text{square}) > 0$. These are the identities
used throughout \texttt{realization/route\_b/}.
\bucket{derived in full (given the standard Brandt row-sum property,
defined with source: Eichler/Pizer)}\chk{D11}
\end{deriv}

\begin{deriv}[D12: the Eichler mass by orbit--stabiliser]\label{d12}
$A_5$ (order 60) acts on $\mathbb{P}^1(\mathbb{F}_q)$ ($q+1$ points);
the ideal classes are the orbits. The mass is
\[
\sum_{\text{orbits}} \frac{1}{|\mathrm{Stab}|}
= \sum_{\text{orbits}} \frac{|\text{orbit}|}{60}
= \frac{q+1}{60},
\]
by orbit--stabiliser and the fact that the orbits partition the
$q + 1$ points. At $q = 31$: orbit sizes $(12, 20)$, stabilisers
$(5, 3)$, mass $\tfrac15 + \tfrac13 = \tfrac{8}{15} = \tfrac{32}{60}$.
The identity $\sum 1/|\mathrm{Stab}| = (q+1)/60$ holds for any $q$
by the same two lines; the orbit sizes quoted are the computed ones at
$q = 31$. The mass check in the realization is this identity --- a
structural sanity gate, automatic once the orbit count is right.
\bucket{derived in full}\chk{D12}
\end{deriv}

\section{Part III: the spectra}

\begin{deriv}[D13: why the tapered prime sum peaks at the
zeros]\label{d13}
The crystal page plots
$F_X(k) = \bigl|\sum_{p \le X} \tfrac{\log p}{\sqrt p}
\bigl(1 - \tfrac{\log p}{\log X}\bigr) p^{ik}\bigr|$. The mechanism:
$F_X$ is (up to prime-power corrections) a Riesz-smoothed partial sum
of the Dirichlet series of $-\zeta'/\zeta$ at $s = \tfrac12 - ik$,
whose analytic continuation has a simple pole at every non-trivial
zero; at $s = \tfrac12 - ik$ with $k > 0$ the nearby pole is the
conjugate zero $\tfrac12 - i\gamma$, so by the conjugation symmetry
of the zeros the peaks appear at $k = |\gamma|$, i.e.\ at the
ordinates. Smoothing converts each nearby pole into a finite peak of
width $\sim 1/\log X$, and the
taper suppresses the Gibbs ringing that a sharp cutoff produces. We
do not derive the full asymptotic here (gap: making the peak shape
and height uniform in $k$ and $X$ is genuine analytic work we have
not done); what we assert is the checkable consequence, and we check
it: every zero in the plotted range owns a local peak, at two
independent prime cutoffs ($X = 10^6$ in the shipped figure, $13/13$;
$X = 2\times10^5$ re-verified independently in the check suite,
$13/13$), with the sharp-cutoff control failing ($5/13$) exactly as
the Gibbs explanation predicts.
\bucket{named as a gap (the uniform peak asymptotic); the mechanism
above is heuristic and its checkable consequence is verified}\chk{D13}
\end{deriv}

\begin{deriv}[D14: the Fibonacci-chain peak module]\label{d14}
The chain is built by the substitution $A \to AB$, $B \to A$ with tile
lengths $\varphi$ and $1$. That its diffraction is pure point with
Bragg peaks on a rank-2 module is the standard cut-and-project theory
of the Fibonacci quasicrystal (defined with source: Baake--Grimm,
\emph{Aperiodic Order} I). The checkable consequence: the strongest
peaks of the computed spectrum lie at
$k = \tfrac{2\pi}{\bar\ell}(m + n\varphi)$ for small integers $m, n$,
where $\bar\ell$ is the mean tile length. The check suite locates the
eight strongest peaks of the shipped 3{,}000-tile spectrum and finds
every one within $0.02$ of that module.
\bucket{defined with source (the consequence checked numerically)}
\chk{D14}
\end{deriv}

\begin{deriv}[D15: the Chebyshev bias mechanism]\label{d15}
Let $\chi$ be the non-trivial character mod 4 and
$\psi(x, \chi) = \sum_{n \le x} \Lambda(n)\chi(n)$,
$\theta(x, \chi) = \sum_{p \le x} (\log p)\chi(p)$. Splitting
$\psi$ over prime powers: the square terms contribute
$\sum_{p^2 \le x} (\log p)\chi(p^2) = \sum_{p \le \sqrt x} \log p
\cdot \chi(p)^2 = \theta(\sqrt x) \sim \sqrt x$, since
$\chi(p)^2 = 1$ for every odd $p$ --- \emph{all prime squares land in
the residue class} $1 \bmod 4$. Cubes and higher contribute
$O(x^{1/3}\log x)$. Hence
\[
\theta(x, \chi) \;=\; \psi(x, \chi) \;-\; \sqrt{x}\,(1 + o(1)),
\]
and since $L(s, \chi)$ has no pole, $\psi(x,\chi)$ is purely a sum of
zero oscillations. Partial summation then gives
\[
\pi(x; 4, 3) - \pi(x; 4, 1)
\;=\; \underbrace{\frac{\sqrt x}{\log x}\,(1 + o(1))}_{\text{the
square-term contribution}} \;+\;
(\text{oscillations from the zeros of } L(s,\chi)),
\]
where the labelled term is the contribution of the prime squares ---
not a dominance asymptotic, since the oscillation term is of the same
order.
This derives the \emph{mechanism} and the scale: the bias coefficient
$D(x)\log x/\sqrt x$ is $O(1)$ (measured: $2.03$, $1.11$, $0.57$ at
$10^6$, $10^7$, $5\times10^7$ --- order one and fluctuating, as the
oscillation term predicts). What it does \emph{not} derive is
dominance: the oscillations are of the same order as the main term, so
``the 3-class leads almost always'' requires the logarithmic-density
analysis of Rubinstein--Sarnak under GRH + linear independence ---
named as a gap, exactly as on the race page.
\bucket{derived in full (the square-term mechanism; dominance is the named gap)}
\chk{D15}
\end{deriv}

\section{The gap ledger, cross-referenced}

For the explorables and explainer, the open items above are exactly
three of the mathematics page's six named gaps: the heuristic status of
the Hardy--Littlewood asymptotic (D7, D8 $\to$ gap 1), the uniform
asymptotics of the tapered spectrum's peaks (D13 $\to$ part of gap 6),
and bias dominance (D15 $\to$ part of gap 6). Everything else used by
the prime pages is either derived in full above or defined with an
explicit source (Ramanujan bound; Sato--Tate law; exact GUE; Fibonacci
cut-and-project; Brandt row-sum property). The verification artifact
\texttt{out/derivations\_check.json} certifies all eighteen checkable
statements.

\begin{thebibliography}{9}\small
\bibitem{HL1923} G.~H. Hardy, J.~E. Littlewood, \emph{Some problems of
`Partitio Numerorum' III}, Acta Math.\ 44 (1923).
\bibitem{BaakeGrimm} M.~Baake, U.~Grimm, \emph{Aperiodic Order,
Vol.~1}, Cambridge Univ.\ Press (2013).
\bibitem{Mehta} M.~L. Mehta, \emph{Random Matrices}, 3rd ed.,
Elsevier (2004).
\bibitem{RubinsteinSarnak} M.~Rubinstein, P.~Sarnak,
\emph{Chebyshev's bias}, Experiment.\ Math.\ 3 (1994).
\bibitem{Eichler} M.~Eichler, \emph{The basis problem for modular
forms and the traces of the Hecke operators}, Springer LNM 320 (1973);
A.~Pizer, \emph{An algorithm for computing modular forms on
$\Gamma_0(N)$}, J.~Algebra 64 (1980).
\bibitem{Cramer} H.~Cram\'er, \emph{On the order of magnitude of the
difference between consecutive prime numbers}, Acta Arith.\ 2 (1936).
\end{thebibliography}

\end{document}
